% Generated draft: requires mathematical and source review.
\documentclass[11pt]{article}
\usepackage[margin=25mm]{geometry}
\usepackage{fontspec}
\setmainfont{Latin Modern Roman}
\usepackage{amsmath,amssymb,mathtools,unicode-math}
\setmathfont{Latin Modern Math}
\usepackage{enumitem,needspace,xurl,hyperref}
\setlength{\parindent}{0pt}
\setlength{\parskip}{5pt}
\setlength{\emergencystretch}{3em}
\newcommand{\N}{\mathbb{N}}
\newcommand{\Z}{\mathbb{Z}}
\newcommand{\Q}{\mathbb{Q}}
\newcommand{\R}{\mathbb{R}}
\newcommand{\C}{\mathbb{C}}
\newcommand{\F}{\mathbb{F}}
\newcommand{\A}{\mathbb{A}}
\newcommand{\PP}{\mathbb{P}}
\newcommand{\E}{\mathbb{E}}
\newcommand{\eps}{\varepsilon}
\newcommand{\dd}{\,\mathrm{d}}
\begin{document}
\section*{UnsolvedMath ID: 10300029}
\textbf{Source problem code:} catalogue entry
\subsection{Definitions and mathematical statement}
A group \(G\) is left-orderable if it admits a strict total order \(<\) such that \(g<h\) implies \(kg<kh\) for every \(k\in G\). Equivalently, there is a positive cone \(P\) with \(PP\subseteq P\) and
\[
G=P\sqcup\{1\}\sqcup P^{-1}.
\]
For a finite symmetric generating set \(S\), let \(B_S(r)\) be the word ball. A customary finite-radius analogue asks for \(P_r\subseteq B_S(r)\setminus\{1\}\) with
\[
B_S(r)\setminus\{1\}=P_r\sqcup P_r^{-1},\qquad
x,y\in P_r,\ xy\in B_S(r)\Longrightarrow xy\in P_r.
\]
Call this predicate \(\mathsf{Ord}_r(G,S)\). The source does not define its finite-radius ordering convention, so this explicit positive-cone version must be checked.

With the intended hyperbolic \(3\)-manifold class and generating-set convention fixed, the first question is whether a single integer \(c\) guarantees
\[
\forall M,\qquad
\mathsf{Ord}_c(\pi_1(M),S_M)\Longrightarrow\pi_1(M)\text{ left-orderable}.
\]
The weaker question asks for an effective procedure that, from the intended finite description of \(M\), returns \(c(M)\) with that implication for \(M\). It need not return the least such radius. A radius computed separately for each manifold is weaker than a universal bound.

Word radius depends on \(S_M\); a rule for choosing generators, or a quantifier over all generating sets, is essential. The source does not state closedness/finite-volume assumptions or an effective input representation. No such manifold or generator restrictions are silently inserted. For the positive-cone convention, global left orders restrict to every finite radius, and compactness relates consistent finite constraints to global orderability; that does not by itself provide a universal or computable detecting radius.

The supplied research summary mentions effective ordering and L-space connections but gives no verified radius result. Its undecidable-in-general/partially-open phrase lacks a precise group class or theorem and is not treated as undecidability for the hyperbolic \(3\)-manifold problem.
\subsection{Short English statement}
Can left-orderability of the intended hyperbolic \(3\)-manifold groups be detected at one universal finite radius, or at a radius effectively computed from the manifold? The local-order test, generators, and input category must be specified.
\subsection{Sources}
\begin{itemize}
\item Local UnsolvedMath registry, permanent ID \(10300029\), exact source code \texttt{catalogue entry}, ``Leaf spaces and transverse structures --- Question 8.6.''
\item Calegari, \emph{Problems in foliations and laminations of 3-manifolds} (2002), Question 8.6; \url{https://arxiv.org/abs/math/0209081}. Source-TeX extraction; recorded access 2026-07-29. The source list contains 67 numbered questions.
\item Supplied OPEN-TRIAGE summary mentions Boyer--Rolfsen--Wiest and Calegari--Dunfield, with L-space/computability connections, but no complete reference or precise theorem resolving either radius question.
\end{itemize}
\end{document}
